
pour tout $x\in X(n)$, et on a
\begin{equation}\tag{1.1.6}
\ell_x=\ell'_x+\ell''_x\qquad (x\in X(n)).
\end{equation}
Alors l'assertion (1.1.3) pour \emph{tous} les membres de (1.1.5) découle immédiatement de l'hypothèse et de l'équation
\begin{equation}\tag{1.1.7}
\cl_\bullet(\underline{G})=\cl_\bullet(\underline{G}')+\cl_\bullet(\underline{G}'').
\end{equation}

On sait que $K_\bullet$ est un foncteur covariant sur la catégorie des schémas localement noethériens par rapport aux morphismes propres. La proposition suivante affirme que notre filtration est fonctorielle :

\textbf{Proposition 1.1.8.} \uline{Soit} $f:X'\to X$ \uline{un morphisme propre de schémas localement noethériens et soit}
\[
f_*:K_\bullet(X')\longrightarrow K_\bullet(X)
\]
\uline{le morphisme induit par} $f$\uline{. Alors}
\begin{equation}\tag{1.1.9}
f_*(\Filt_n(X'))\subset\Filt_n(X)
\end{equation}
\uline{pour tout entier} $n$.

\emph{Démonstration :} Si $\underline{F}$ est un $\underline{O}_{X'}$-module cohérent, alors on a par définition (IV 2.11) :
\[
f_*(\cl_\bullet(\underline{F}))=\sum_{i\in\mathbb N}(-1)^i\cl_\bullet(R^if_*(\underline{F})).
\]
Or, les $\underline{O}_X$-modules $R^if_*(\underline{F})$ sont concentrés sur $f(\Supp(\underline{F}))$ et d'après (EGA IV 5.4.1 (ii)) on a
\[
\dim(\Supp(\underline{F}))\geq\dim(f(\Supp(\underline{F}))),
\]
d'où l'assertion.

